\documentclass[12pt]{article}
\usepackage[utf8]{inputenc}
\usepackage[T1]{fontenc}
\usepackage[margin=1in]{geometry}
\usepackage{setspace}
\usepackage{xcolor}
\usepackage[colorlinks=true, urlcolor=blue!60!black, linkcolor=blue!60!black]{hyperref}
\setlength{\parindent}{2em}
\setlength{\parskip}{1.25em}
\renewcommand{\baselinestretch}{1.0}
\title{\textit{Plyler v. Doe}, 457 U.S. 202 (1982)}
\author{Case Brief}
\date{}
\begin{document}
\maketitle

% Case Brief: Plyler v. Doe, 457 U.S. 202 (1982)
% Constitutional Law Briefing Book — Chapter: Equal Protection

\section{Case Metadata}

\begin{itemize}
    \item \textbf{Case Name:} \textit{Plyler, Superintendent, Tyler Independent School District, et al.\ v.\ Doe et al.}, consolidated with \textit{Texas et al.\ v.\ Certain Named and Unnamed Undocumented Alien Children et al.}
    \item \textbf{Citation:} 457 U.S. 202, 102 S. Ct. 2382, 72 L. Ed. 2d 786 (1982) (\href{https://www.law.cornell.edu/supremecourt/text/457/202}{Cornell LII}; \href{https://supreme.justia.com/cases/federal/us/457/202/}{Justia})
    \item \textbf{Docket Numbers:} Nos.\ 80-1538 and 80-1934
    \item \textbf{Dates:} Argued December 1, 1981; Decided June 15, 1982; Rehearings denied September 9, 1982 (458 U.S. 1131)
    \item \textbf{Vote:} 5--4
    \item \textbf{Majority Opinion:} Justice William J. Brennan, Jr., joined by Justices Thurgood Marshall, Harry A. Blackmun, Lewis F. Powell, Jr., and John Paul Stevens
    \item \textbf{Concurring Opinions:} Justice Marshall (concurring); Justice Blackmun (concurring); Justice Powell (concurring)
    \item \textbf{Dissenting Opinion:} Chief Justice Warren E. Burger, joined by Justices Byron R. White, William H. Rehnquist, and Sandra Day O'Connor
    \item \textbf{Constitutional Provision and Statute at Issue:} Equal Protection Clause of the Fourteenth Amendment (``nor deny to any person within its jurisdiction the equal protection of the laws''); Texas Education Code \S~21.031 (Vernon Supp.\ 1981)
\end{itemize}

\section{Facts and Procedural History}

In May 1975, the Texas Legislature revised \href{https://www.law.cornell.edu/supremecourt/text/457/202}{Texas Education Code \S~21.031} to withhold state educational funds from local school districts for the education of children who were not ``legally admitted'' into the United States, and to authorize local school districts to deny free enrollment in public schools to such children. Specifically, \S~21.031(a)--(c) restricted entitlement to the Available School Fund and tuition-free public school admission to children between ages five and twenty-one who were either ``citizens of the United States or legally admitted aliens.''

Two consolidated lines of litigation challenged \S~21.031:

\begin{enumerate}
    \item \textbf{\textit{Plyler v.\ Doe} (No.\ 80-1538):} Despite the 1975 amendment, the Tyler Independent School District (ISD) in Smith County, Texas, continued enrolling undocumented children free of charge until the 1977--1978 school year. In July 1977, Tyler ISD adopted a policy requiring undocumented children to pay a ``full tuition fee'' (\$1,000 per year) to enroll, defining a ``legally admitted alien'' as someone with federal documentation of lawful presence or confirmation from the Immigration and Naturalization Service (INS) that documentation was being processed. In September 1977, undocumented school-age children of Mexican origin residing in Smith County filed a class action in the U.S. District Court for the Eastern District of Texas against Superintendent James Plyler and the Board of Trustees; the State of Texas intervened as a defendant. Following an extensive evidentiary hearing, the District Court (\textit{Doe v.\ Plyler}, 458 F. Supp. 569 (E.D. Tex. 1978)) found that \S~21.031 neither had the purpose nor the effect of deterring illegal immigration, that excluding undocumented children would not necessarily improve educational quality, and that many undocumented children would remain permanently in the United States. The court permanently enjoined the policy under both the Equal Protection Clause (finding it failed even rational basis review) and the Supremacy Clause (federal preemption). The Fifth Circuit (\textit{Doe v.\ Plyler}, 628 F.2d 448 (5th Cir. 1980)) reversed on federal preemption but affirmed on Equal Protection grounds, holding \S~21.031 unconstitutional under either rational basis or heightened scrutiny.
    \item \textbf{\textit{In re Alien Children Education Litigation} (No.\ 80-1934):} During 1978 and 1979, additional suits challenging \S~21.031 were filed in the Southern, Western, and Northern Districts of Texas against the State of Texas and the Texas Education Agency. In November 1979, the Judicial Panel on Multidistrict Litigation consolidated the statewide claims in the U.S. District Court for the Southern District of Texas. In July 1980, that court (\textit{In re Alien Children Education Litigation}, 501 F. Supp. 544 (S.D. Tex. 1980)) rejected federal preemption but held that the ``absolute deprivation of education'' triggered strict judicial scrutiny and invalidated \S~21.031 under the Equal Protection Clause. On February 23, 1981, the Fifth Circuit summarily affirmed on the authority of its decision in \textit{Plyler}.
\end{enumerate}

The Supreme Court noted probable jurisdiction in both cases (451 U.S. 968 (1981); 452 U.S. 937 (1981)) and consolidated them for briefing and oral argument. Because the Court resolved the dispute under the Fourteenth Amendment, it did not reach the plaintiffs' renewed Supremacy Clause preemption claims (\href{https://www.law.cornell.edu/supremecourt/text/457/202}{457 U.S. at 210 n.8}).

\section{Issue(s) Presented}

\begin{enumerate}
    \item Are undocumented immigrants physically present inside a state's borders ``persons within its jurisdiction'' entitled to the protection of the Fourteenth Amendment's Equal Protection Clause?
    \item Consistent with the Equal Protection Clause, may a state withhold state funds from local school districts for the education of undocumented school-age children and authorize local districts to deny them the free public K--12 education provided to citizens and legally admitted alien children?
\end{enumerate}

\section{Holding}

\begin{enumerate}
    \item \textbf{Yes} (unanimous on the threshold jurisdictional question). Undocumented immigrants physically present within a state's territorial boundaries are ``persons within its jurisdiction'' protected by the Equal Protection Clause of the Fourteenth Amendment (\href{https://www.law.cornell.edu/supremecourt/text/457/202}{457 U.S. at 210--16}; \textit{id.}\ at 243 (Burger, C.J., dissenting)).
    \item \textbf{No} (5--4). Although undocumented aliens are not a ``suspect class'' and public education is not a ``fundamental right'' triggering strict scrutiny, Texas Education Code \S~21.031 imposes a lifetime hardship on a discrete class of innocent children not accountable for their disabling status. Under a heightened rational basis (intermediate-like) standard, the discrimination cannot be sustained unless it furthers some \textbf{substantial goal of the State}. Because Texas failed to show that excluding undocumented children furthered any substantial state interest, \S~21.031 violates the Equal Protection Clause (\href{https://www.law.cornell.edu/supremecourt/text/457/202}{457 U.S. at 223--30}).
\end{enumerate}

\section{Reasoning}

Justice Brennan's majority opinion proceeded in four analytical steps:

\subsection*{1. Territorial Scope of ``Within Its Jurisdiction''}
Texas argued at the threshold that undocumented immigrants, because their entry violated federal law, were not ``within its jurisdiction'' under the Equal Protection Clause. The Court flatly rejected this reading. First, an alien---regardless of immigration status---is undeniably a ``person'' in the ordinary sense, as long recognized under the Fifth and Fourteenth Amendment Due Process Clauses (\textit{Yick Wo v.\ Hopkins}, 118 U.S. 356, 369 (1886); \textit{Wong Wing v.\ United States}, 163 U.S. 228, 238 (1896); \textit{Mathews v.\ Diaz}, 426 U.S. 67, 77 (1976)). Second, the Equal Protection Clause and Due Process Clause protect an identical class of persons. Reviewing the Reconstruction debates---including statements by Representative John Bingham and Senator Jacob Howard---the Court concluded that ``within its jurisdiction'' was designed in a territorial sense to guarantee equal protection to ``anyone, citizen or stranger, who is subject to the laws of a State'' (\href{https://www.law.cornell.edu/supremecourt/text/457/202}{457 U.S. at 215}). Allowing a state to define persons physically subject to its civil and criminal laws as outside its jurisdiction would invert the Clause's core anti-caste purpose.

\subsection*{2. Selecting the Standard of Review: Heightened Rational Basis / Intermediate-Like Scrutiny}
The Court declined to apply strict scrutiny, reaffirming two established doctrinal boundaries:
\begin{itemize}
    \item \textbf{Not a Suspect Class:} Undocumented aliens are not a suspect class because entry into the class is the product of voluntary, unlawful action and is not a ``constitutional irrelevancy'' (\href{https://www.law.cornell.edu/supremecourt/text/457/202}{457 U.S. at 219 n.19, 223}).
    \item \textbf{Not a Fundamental Right:} Public education is not a right explicitly or implicitly guaranteed by the Constitution (\href{https://supreme.justia.com/cases/federal/us/411/1/}{\textit{San Antonio Independent School District v.\ Rodriguez}}, 411 U.S. 1, 35 (1973)).
\end{itemize}

Yet the Court refused to apply ordinary, deferential rational basis review because of the convergence of two extraordinary factors:
\begin{itemize}
    \item \textbf{Innocent Children and Anti-Caste Principles:} Unlike adult unlawful entrants, minor children ``can affect neither their parents' conduct nor their own status'' (\textit{Trimble v.\ Gordon}, 430 U.S. 762, 770 (1977)). Drawing on illegitimacy precedents (\textit{Weber v.\ Aetna Casualty \& Surety Co.}, 406 U.S. 164, 175 (1972)), the Court reasoned that visiting condemnation upon children for their parents' misconduct is contrary to basic principles of individual responsibility and risks creating a permanent ``shadow population'' or underclass incompatible with the Fourteenth Amendment.
    \item \textbf{The Singular Importance of Basic Education:} While not a fundamental right, public education is not merely an ordinary social-welfare benefit. Citing \textit{Brown v.\ Board of Education}, 347 U.S. 483, 493 (1954), \textit{Wisconsin v.\ Yoder}, 406 U.S. 205, 221 (1972), and \textit{Meyer v.\ Nebraska}, 262 U.S. 390, 400 (1923), the Court emphasized that education sustains democratic institutions, transmits cultural values, and equips individuals for economic self-sufficiency. Complete denial of basic education inflicts the lifelong stigma of illiteracy.
\end{itemize}

Combining these two factors, the Court held that \S~21.031 ``can hardly be considered rational unless it furthers some substantial goal of the State'' (\href{https://www.law.cornell.edu/supremecourt/text/457/202}{457 U.S. at 224})---an intermediate-like standard demanding a substantial state interest rather than mere conceivable rationality.

\subsection*{3. Absence of Congressional Authorization}
Texas contended that federal disapproval of unlawful entry supplied a rational basis for state exclusion. The Court rejected this argument, observing that while Congress holds plenary authority over immigration and naturalization (U.S. Const.\ art.\ I, \S~8, cl.\ 4), states enjoy no independent power to classify aliens (\textit{Hines v.\ Davidowitz}, 312 U.S. 52 (1941)). Unlike the state employment restriction upheld in \textit{DeCanas v.\ Bica}, 424 U.S. 351 (1976), which mirrored congressional policy protecting domestic labor, no federal policy supported denying elementary education to undocumented children. Moreover, because the federal government exercises discretionary relief from deportation, a state cannot predict which undocumented children will remain permanently or eventually adjust to lawful status.

\subsection*{4. Failure of Texas's Asserted State Interests}
Evaluating Texas's overarching claim of preserving limited fiscal resources---which standing alone cannot justify a discriminatory classification (\textit{Graham v.\ Richardson}, 403 U.S. 365, 374--75 (1971))---the Court rejected three specific state interests:
\begin{enumerate}
    \item \textbf{Deterring Illegal Immigration:} The trial record showed that undocumented immigrants generally contributed labor and taxes while underutilizing public services, and that the dominant magnet for unlawful entry was employment, not free schooling. Charging tuition to children was a ``ludicrously ineffectual attempt to stem the tide of illegal immigration'' compared to restricting unlawful employment (\href{https://www.law.cornell.edu/supremecourt/text/457/202}{457 U.S. at 228}).
    \item \textbf{Preserving Educational Quality:} Texas produced no credible evidence that excluding undocumented children would improve the quality of public education, particularly since state and federal funding scaled with enrollment and undocumented children were ``basically indistinguishable'' in educational cost and need from legally resident alien children (\textit{id.}\ at 229).
    \item \textbf{Lower Likelihood of Remaining in Texas:} The State could not guarantee that citizen or lawful-resident children would remain in Texas after graduation, whereas the record showed many undocumented children would stay indefinitely. Promoting a permanent ``subclass of illiterates'' would impose far greater long-term societal costs (unemployment, welfare, and crime) than any short-term fiscal savings (\textit{id.}\ at 230).
\end{enumerate}

The Court also dismissed Tyler ISD's oral-argument attempt to recharacterize \S~21.031 as a neutral residency rule, noting that a state may not circumvent the Equal Protection Clause simply by defining a disfavored group as ``nonresident'' (\textit{id.}\ at 227 n.22).

\section{Concurring and Dissenting Opinions}

\begin{itemize}
    \item \textbf{Justice Marshall, Concurring (\href{https://www.law.cornell.edu/supremecourt/text/457/202}{457 U.S. at 230--31}):} Joined the majority opinion while reiterating his dissent in \href{https://supreme.justia.com/cases/federal/us/411/1/}{\textit{San Antonio Independent School District v.\ Rodriguez}}, 411 U.S. 1, 70--133 (1973). He maintained that education is a fundamental interest and that \textit{Plyler} vindicated his flexible ``sliding-scale'' approach to Equal Protection, calibrating scrutiny according to both the societal importance of the interest affected and the invidiousness of the classification.
    \item \textbf{Justice Blackmun, Concurring (\href{https://www.law.cornell.edu/supremecourt/text/457/202}{457 U.S. at 231--36}):} Having joined the majority in \textit{Rodriguez}, Justice Blackmun wrote separately to reconcile the two cases. He analogized the \textit{complete denial} of basic education to the denial of the right to vote: though neither is explicitly enumerated in the Constitution, denying either relegates an identifiable group to permanent second-class status, striking at the heart of Equal Protection. Unlike the funding disparities in \textit{Rodriguez}, \S~21.031 effected an absolute deprivation. He also stressed federal preemption concerns: because states cannot determine which undocumented children will ultimately be permitted to remain, \S~21.031 was fatally imprecise.
    \item \textbf{Justice Powell, Concurring (\href{https://www.law.cornell.edu/supremecourt/text/457/202}{457 U.S. at 236--42}):} Author of the \textit{Rodriguez} majority opinion, Justice Powell emphasized the ``unique character'' of \textit{Plyler}. Analogizing undocumented children to illegitimate children under \textit{Weber}, \textit{Trimble}, and \textit{Lalli v.\ Lalli}, 439 U.S. 259 (1978), he endorsed heightened intermediate scrutiny requiring that the state's means bear a ``fair and substantial relation'' to substantial state interests. He distinguished \textit{Rodriguez} because \S~21.031 singled out a discrete class of innocent children for total exclusion based on parental misconduct, while expressing sympathy for border states burdened by Congress's lack of ``effective leadership'' on immigration.
    \item \textbf{Chief Justice Burger, Dissenting, joined by Justices White, Rehnquist, and O'Connor (\href{https://www.law.cornell.edu/supremecourt/text/457/202}{457 U.S. at 242--54}):} Chief Justice Burger agreed that undocumented aliens are ``within the jurisdiction'' of Texas and candidly remarked that, as a matter of social policy, depriving children of an elementary education is ``senseless'' and ``folly.'' Nevertheless, he accused the majority of acting as ``Platonic Guardians'' and engaging in unabashedly result-oriented policymaking to compensate for congressional inaction. Once the majority conceded that undocumented aliens are not a suspect class and education is not a fundamental right, settled doctrine mandated ordinary rational basis review. Under that deferential standard, Texas had a legitimate fiscal interest in conserving finite revenues for lawful residents, reinforced by federal exclusions of undocumented aliens from welfare programs (such as Food Stamps, AFDC, and Medicare). He criticized the majority's ``quasi-suspect-class and quasi-fundamental-rights'' patchwork as unprincipled, warned that \textit{Plyler} could not logically distinguish education from food, shelter, or health care, and predicted that the ruling's ``unique confluence of theories and rationales'' would ``likely stand for little beyond the results in these particular cases.''
\end{itemize}

\section{Key Quotations}

\begin{enumerate}
    \item \textbf{Territorial Reach of the Equal Protection Clause (Brennan, J., Majority):}
    \begin{quote}
        ``Use of the phrase `within its jurisdiction' thus does not detract from, but rather confirms, the understanding that the protection of the Fourteenth Amendment extends to anyone, citizen or stranger, who is subject to the laws of a State, and reaches into every corner of a State's territory.''
    \end{quote}
    --- \href{https://www.law.cornell.edu/supremecourt/text/457/202}{457 U.S. at 215} (\href{https://supreme.justia.com/cases/federal/us/457/202/}{Justia}).

    \item \textbf{Special Constitutional Status of Public Education (Brennan, J., Majority):}
    \begin{quote}
        ``Public education is not a `right' granted to individuals by the Constitution. \textit{San Antonio Independent School Dist.\ v.\ Rodriguez}, 411 U.S. 1, 35 (1973). But neither is it merely some governmental `benefit' indistinguishable from other forms of social welfare legislation. Both the importance of education in maintaining our basic institutions, and the lasting impact of its deprivation on the life of the child, mark the distinction.''
    \end{quote}
    --- \href{https://www.law.cornell.edu/supremecourt/text/457/202}{457 U.S. at 221} (\href{https://supreme.justia.com/cases/federal/us/457/202/}{Justia}).

    \item \textbf{Heightened Scrutiny for Lifetime Hardship on Innocent Children (Brennan, J., Majority):}
    \begin{quote}
        ``Section 21.031 imposes a lifetime hardship on a discrete class of children not accountable for their disabling status. The stigma of illiteracy will mark them for the rest of their lives. By denying these children a basic education, we deny them the ability to live within the structure of our civic institutions, and foreclose any realistic possibility that they will contribute in even the smallest way to the progress of our Nation. In determining the rationality of \S~21.031, we may appropriately take into account its costs to the Nation and to the innocent children who are its victims. In light of these countervailing costs, the discrimination contained in \S~21.031 can hardly be considered rational unless it furthers some substantial goal of the State.''
    \end{quote}
    --- \href{https://www.law.cornell.edu/supremecourt/text/457/202}{457 U.S. at 223--24} (\href{https://supreme.justia.com/cases/federal/us/457/202/}{Justia}).

    \item \textbf{Rejection of Fiscal Savings Against Creation of an Illiterate Subclass (Brennan, J., Majority):}
    \begin{quote}
        ``It is difficult to understand precisely what the State hopes to achieve by promoting the creation and perpetuation of a subclass of illiterates within our boundaries, surely adding to the problems and costs of unemployment, welfare, and crime. It is thus clear that whatever savings might be achieved by denying these children an education, they are wholly insubstantial in light of the costs involved to these children, the State, and the Nation.''
    \end{quote}
    --- \href{https://www.law.cornell.edu/supremecourt/text/457/202}{457 U.S. at 230} (\href{https://supreme.justia.com/cases/federal/us/457/202/}{Justia}).

    \item \textbf{Judicial Restraint and the ``Platonic Guardians'' Critique (Burger, C.J., Dissenting):}
    \begin{quote}
        ``Were it our business to set the Nation's social policy, I would agree without hesitation that it is senseless for an enlightened society to deprive any children---including illegal aliens---of an elementary education. \ldots\ However, the Constitution does not constitute us as `Platonic Guardians' nor does it vest in this Court the authority to strike down laws because they do not meet our standards of desirable social policy, `wisdom,' or `common sense.' ''
    \end{quote}
    --- \href{https://www.law.cornell.edu/supremecourt/text/457/202}{457 U.S. at 242} (\href{https://supreme.justia.com/cases/federal/us/457/202/}{Justia}).
\end{enumerate}

\section{Doctrinal Significance and Modern Legacy (Through 2026)}

\subsection*{1. Doctrinal Containment and ``Rational Basis with Bite''}
Chief Justice Burger's prediction that \textit{Plyler} would stand as a narrow, fact-bound outlier proved largely accurate. Within one year, \href{https://supreme.justia.com/cases/federal/us/461/321/}{\textit{Martinez v.\ Bynum}}, 461 U.S. 321 (1983), upheld a separate portion of Texas Education Code \S~21.031 imposing a bona fide residency requirement for public schooling. In \href{https://supreme.justia.com/cases/federal/us/487/450/}{\textit{Kadrmas v.\ Dickinson Public Schools}}, 487 U.S. 450, 459 (1988), the Court refused to extend \textit{Plyler}'s heightened review to school-bus user fees, stressing that \textit{Plyler} rested on ``unique circumstances.'' Under PRWORA and IIRIRA (1996), lower courts have likewise confined \textit{Plyler} to K--12 education, generally upholding state bars on undocumented students' access to public higher education, in-state tuition, and financial aid (\href{https://www.everycrsreport.com/reports/R43447.html}{CRS Report R43447}). Still, \textit{Plyler} endures as a foundational ``rational basis with bite'' precedent alongside \textit{Cleburne} and \textit{Romer}.

\subsection*{2. Practical Impact and Renewed Challenges (2025--2026)}
A December 2025 analysis estimated that at least 4.8 million students have benefited from \textit{Plyler} since 1982 (\href{https://www.k12dive.com/news/what-would-happen-if-undocumented-children-couldnt-attend-public-schools/807318/}{K--12 Dive}; \href{https://www.americanimmigrationcouncil.org/fact-sheet/plyler-v-doe-public-education-noncitizen-students/}{American Immigration Council, Aug.\ 2026}). Federal courts previously blocked state barriers including California's Proposition 187 and Alabama's H.B.\ 56 (2011) (\href{https://www.theguardian.com/us-news/2025/jan/28/oklahoma-schools-immigration-status}{The Guardian}). In 2025--2026, however, lawmakers in at least five states---including Tennessee, Oklahoma, Texas, Indiana, and New Jersey (e.g., A.B.\ 323)---pursued measures to collect student immigration status, charge tuition, or deny enrollment (\href{https://bellwether.org/publications/equal-protection-and-public-education-access-for-immigrant-students-a-50-state-review/}{Bellwether, May 2026}; \href{https://www.americanimmigrationcouncil.org/fact-sheet/plyler-v-doe-public-education-noncitizen-students/}{American Immigration Council, Aug.\ 2026}). Tennessee's Senate passed S.B.\ 836 in April 2025 and its House passed amended H.B.\ 793 in March 2026, but both bills died upon adjournment on April 23, 2026 (\href{https://ace-usa.org/blog/research/research-immigration/should-public-schools-collect-students-immigration-status-a-look-at-tennessees-debate/}{ACE-USA, Oct.\ 2026}). Oklahoma's State Board of Education similarly proposed a 2025 enrollment-verification rule (\href{https://www.theguardian.com/us-news/2025/jan/28/oklahoma-schools-immigration-status}{The Guardian}). Following a February 2026 \href{https://www.heritage.org/border-security/report/every-state-should-challenge-plyler-v-doe-time-end-free-education-illegal-0}{Heritage Foundation report} urging states to test \textit{Plyler}, the U.S. House Judiciary Committee held a March 18, 2026 hearing, \href{https://judiciary.house.gov/committee-activity/hearings/immigration-policy-court-order-adverse-effects-plyler-v-doe}{\textit{Immigration Policy by Court Order: The Adverse Effects of Plyler v.\ Doe}}. As of late 2026, \textit{Plyler} remains binding nationwide (\href{https://bellwether.org/publications/equal-protection-and-public-education-access-for-immigrant-students-a-50-state-review/}{Bellwether, May 2026}).

\section{Pedagogical Critiques and Exam Hypotheticals}

\subsection*{1. Doctrinal Critiques and the ``Opposite Side'' Contrast}
Critics fault \textit{Plyler} for combining two non-triggers (a non-suspect class and a non-fundamental right) into ad hoc intermediate scrutiny, and argue its logic fits Supremacy Clause preemption better than Equal Protection. For exam mastery, contrast \textit{Plyler} with its ``Opposite Side'' precedents upholding educational disparities under rational basis review:
\begin{itemize}
    \item \textbf{\href{https://supreme.justia.com/cases/federal/us/411/1/}{\textit{San Antonio ISD v.\ Rodriguez}} (1973):} Upheld Texas property-tax funding disparities because every child still received a basic education; \textit{Plyler} struck down an \textit{absolute, status-based deprivation} targeting innocent children.
    \item \textbf{\href{https://supreme.justia.com/cases/federal/us/487/450/}{\textit{Kadrmas v.\ Dickinson Public Schools}} (1988):} Upheld school-bus fees because students were neither penalized for parental unlawful status nor legally barred from school.
\end{itemize}

\subsection*{2. Exam Hypotheticals}
\begin{enumerate}
    \item \textbf{Status Collection Without Tuition:} State X requires schools to verify and report aggregate student immigration status while keeping enrollment free. \textit{Issue:} Whether the chilling effect on enrollment functions as an unconstitutional barrier under \textit{Plyler} or survives rational basis review as fiscal data collection under \textit{Kadrmas}.
    \item \textbf{Higher Education and Unaccompanied Minors:} State Y denies in-state university tuition to undocumented graduates and charges K--12 tuition to 17-year-old unaccompanied minors. \textit{Issue:} The university ban likely survives (\href{https://www.everycrsreport.com/reports/R43447.html}{CRS Report R43447}), whereas the K--12 rule squarely tests \href{https://www.law.cornell.edu/supremecourt/text/457/202}{Justice Powell's concurrence (457 U.S. at 240 n.5)}, which reserved the case of a mature minor entering unlawfully on his own volition.
    \item \textbf{Congressional Authorization:} Congress expressly authorizes states to deny free K--12 schooling to undocumented children. \textit{Issue:} Whether federal authorization supplies the missing substantial interest under Part IV of \textit{Plyler}, or whether Congress cannot authorize states to violate the Fourteenth Amendment.
\end{enumerate}

\end{document}
